%% Dans ce fichier, les �ventuels accents sont cod�s
%% en ISO-8859-15 (aka latin9)
%% La ligne suivante donne trois exemples 
%% e accent aigu : �, e accent grave : �, e-dans-l'o : �
%%%%%ENTREES%%%%
\begin{tabarg}
\multicolumn{3}{l}{\titretab {Entr\'ees}}\\
\hline
$ X $ & Vecteur, \texttt{float} & Param�tre � traiter (unit� quelconque)\\
$N_{bpt}$ & Coeff., \texttt{int} & Nombre de points sur lequel on calcule la moyenne glissante\\
\end{tabarg}

\medskip
%%%%%SORTIES%%%%
\begin{tabarg}
\multicolumn{3}{l}{\titretab {Sorties}}\\
\hline
$X_f$ & Vecteur, \texttt{float} & Param�tre r�sultat (m�me unit� que $X$)\\
$Flag$ & Vecteur, \texttt{float} & Flag indiquant les valeurs remplac�es (0 ou 1)\\
\end{tabarg}
